\section{Verification environment}
\label{sec:env}

The testbench reuses the HW4 checks, split into layers that each do one job.
The driver is the only code that sets DUT inputs, and it drives on the falling
edge. The monitor only reads pins and records the inputs the DUT received and
its outputs. The scoreboard is a reference queue that predicts the count, the
flags, and the next word out. \texttt{FifoEnv} runs one clock cycle through
all three and records coverage. Directed and random tests record coverage
under separate names, so directed cycles never count toward random closure.

\paragraph{Checking every cycle.} The monitor samples just before and just
after each rising edge, and the scoreboard checks both. A wrong count or flag
(F5) or output word (F6) is caught on the cycle it appears. Reset uses the
same path and must leave the FIFO empty (F1).

\paragraph{Guarding against a broken driver.} The scoreboard updates from the
inputs the monitor saw, so a driver bug could fool it. Each cycle therefore
first checks that the sampled inputs match the request, and names the pin on
a mismatch.

\paragraph{Validation.} After the split, the clean run passed at the same
simulated time, coverage was identical to the earlier baseline, and the
released pool stayed 8/8. Three bugs added on purpose (the monitor sampling
too early, the driver never asserting read, the driver returning too early)
were each caught. Logs are in \path{evidence/driver-monitor/}.

\paragraph{Limits.} A glitch between the two samples is missed. Only an
8-bit, 8-deep FIFO was tested.
